% Academic CV of Vincenzo Di Florio.
% Build: cd cv && make   (writes ../docs/assets/pdf/CV.pdf)
% Look and feel: cv-style.sty (same palette and typefaces as vdiflorio.github.io).
\documentclass[11pt]{article}
\usepackage{cv-style}
\hypersetup{pdftitle={Vincenzo Di Florio — Curriculum vitae}, pdfauthor={Vincenzo Di Florio}}

\begin{document}

\cvheader{Curriculum vitae}

%==============================================================================
\section{Current position}
\entry{Post-doctoral Research Fellow, MOX – Modeling and Scientific Computing}{Jan 2025 – present}{Department of Mathematics, Politecnico di Milano, Milan, Italy}

%==============================================================================
\section{Research interests}
My research focuses on:
\begin{itemize}
  \item protein–protein binding through continuum electrostatics, with a focus on electrostatic complementarity at biomolecular interfaces, and efficient numerical solvers for the Poisson–Boltzmann equation. I also develop data-driven and physics-guided neural surrogates that accelerate electrostatic calculations;
  \item non-equilibrium statistical dynamics and the behaviour of low-dimensional systems.
\end{itemize}

%==============================================================================
\section{Education}
\entry{Ph.D. in Pure and Applied Mathematics [EQF 8]}{2021 – 2025}{Politecnico di Torino, Turin, Italy, and Istituto Italiano di Tecnologia, Genoa, Italy}
\detail{Research topics}{biomolecular modelling, continuum electrostatics, and the development and customisation of protocols and software tools for biological systems}
\detail{Thesis}{Mathematical, Algorithmic and Numerical Solutions to Enhance Electrostatic Calculations for Biomolecules in Electrolytic Solutions}
\detail{Supervisors}{Dr. Walter Rocchia, Prof. Lamberto Rondoni}

\entry{Master's degree in Physics of Complex Systems [EQF 7]}{2019 – 2021}{Università degli Studi di Torino, Turin, Italy}
\detail{Grade}{110/110 cum laude, with distinction from the examination board}
\detail{Related coursework}{statistical physics and mathematical modelling}
\detail{Thesis}{State equations and order fluctuations in 1D and 3D}
\detail{Advisor}{Prof. Lamberto Rondoni}

\entry{Bachelor's degree in Physics [EQF 6]}{2016 – 2019}{Università degli Studi di Trieste, Trieste, Italy}
\detail{Grade}{110/110 cum laude}
\detail{Related coursework}{major in physics, with minor courses in numerical analysis and dynamical systems}
\detail{Thesis}{Quantum dots for quantum computation}
\detail{Advisor}{Prof. Angelo Bassi}

%==============================================================================
\section{Research experience}
\entry{Post-doctoral Research Fellow, Politecnico di Milano}{Jan 2025 – present}{Milan, Italy}
\begin{itemize}
  \item Advanced models and efficient numerical methods for molecular electrostatics.
\end{itemize}

\entry{Research collaboration, CONCEPT Lab, Istituto Italiano di Tecnologia}{Nov 2024 – Dec 2024}{Genoa, Italy}
\begin{itemize}
  \item Selected the theoretical model and assessed the feasibility of a code for the nonlinear Poisson–Boltzmann equation.
\end{itemize}

\entry{Internship, City College of New York (CUNY)}{May 2024 – Aug 2024}{New York, USA}
\begin{itemize}
  \item Integrated NextGenPB as the electrostatic solver in the MCCE workflow.
  \item Developed a benchmark application for the calculation of p\textit{K}\textsubscript{a} values.
\end{itemize}

\entry{Internship, SISSA}{Sep 2018 – Dec 2018}{Trieste, Italy}
\begin{itemize}
  \item Sparse identification of dynamical systems from data.
\end{itemize}

%==============================================================================
\section{Publications}
\subsection{Journal articles}
\begin{itemize}
  \item Alharbi, S., Kumawat, A., Huntington, B., Momin, A. A., Alfaiyz, S., \textbf{Di Florio, V.}, Rigoli, M., Cardona Londono, K., Rocchia, W., Sandholu, A., \& Arold, S. T. (2026). \emph{SCRIB inhibits PP1 phosphatase activity without blocking or structurally altering its active site.} Journal of Biological Chemistry, 302(11), 113599. \doi{10.1016/j.jbc.2026.113599}
  \item Carbone, D., \textbf{Di Florio, V.}, Lepri, S., \& Rondoni, L. (2026). \emph{Computing nonequilibrium transport from short-time transients: from Lorentz gas to heat conduction in one-dimensional chains.} The Journal of Chemical Physics, 164(10), 104108. \doi{10.1063/5.0320325}
  \item Mazzeo, M. D., \textbf{Di Florio, V.}, Rocchia, W., \& Decherchi, S. (2026). \emph{Patch-based ray tracing in NanoShaper boosts molecular surface computation.} Journal of Chemical Information and Modeling, 66(1), 461–471. \doi{10.1021/acs.jcim.5c02287}
  \item \textbf{Di Florio, V.}, Ansalone, P., Siryk, S. V., Decherchi, S., De Falco, C., \& Rocchia, W. (2025). \emph{NextGenPB: an analytically-enabled super resolution tool for solving the Poisson–Boltzmann equation featuring local (de)refinement.} Computer Physics Communications, 317, 109816. \doi{10.1016/j.cpc.2025.109816}
  \item \textbf{Di Florio, V.}, Giberti, C., Rondoni, L., \& Zhao, H. (2024). \emph{Microscopic state equation for oscillator chains.} The European Physical Journal Plus, 139(7), 622. \doi{10.1140/epjp/s13360-024-05419-1}
  \item Rondoni, L., \& \textbf{Di Florio, V.} (2024). \emph{Probability turns material: the Boltzmann equation.} Entropy, 26(2), 171. \doi{10.3390/e26020171}
\end{itemize}

%==============================================================================
\section{Presentations}
\subsection{Contributed talks}
\begin{itemize}
  \item “Data-driven resolution enhancement of electrostatic potentials using NextGenPB”, GIMC SIMAI Young 2026, Pisa, Italy (Jun 2026)
  \item “NextGenPB: a modern parallel finite element solver for the Poisson–Boltzmann equation”, SIAM Conference on Parallel Processing for Scientific Computing (PP26), Berlin, Germany (Mar 2026)
  \item “State equation for oscillator chain”, ICIAM 2023, Tokyo, Japan (Aug 2023)
\end{itemize}
\subsection{Posters}
\begin{itemize}
  \item “Mapping electrostatic complementarity at biomolecular interfaces with NextGenPB”, EuroQSAR 2026, Perugia, Italy (Sep–Oct 2026)
  \item “Biomolecular electrostatics applications of the NextGenPB solver”, Protein Electrostatics 2025, Lisbon, Portugal (Jun 2025)
  \item “NextGenPB: a novel Poisson–Boltzmann solver leveraging super-resolution and local refinement”, ACS Fall 2024, Denver, USA (Aug 2024)
  \item “Efficient and scalable implementation of a linearized Poisson–Boltzmann solver on hierarchically refined Cartesian meshes”, Protein Electrostatics 2023, Genoa, Italy (Jun 2023)
\end{itemize}

%==============================================================================
\section{Awards and grants}
\begin{itemize}
  \item \textbf{PRACE-FENIX 2025 Call for Proposals}, Principal Investigator (proposal FENIX2501-050), Sep 2026 – Sep 2027: computing resources awarded at GAIA Cloud (64 vCPU and 30\,000 GB of storage, plus 16 vCPU and 2\,000 GB of storage) for 12 months, after evaluation of 14 proposals.
\end{itemize}

%==============================================================================
\section{Teaching}
\entry{Teaching assistant – Advanced Methods for Scientific Computing (AMSC)}{2026 – present}{Department of Mathematics, Politecnico di Milano, Milan, Italy}
\entry{Tutor – Mathematical Physics (Fisica Matematica)}{2026 – present}{School of Industrial and Information Engineering, Politecnico di Milano, Milan, Italy}
\entry{Tutor – Methods for Environmental Engineering}{Feb 2024 – Apr 2024}{Politecnico di Torino, Turin, Italy}
\entry{Teaching assistant – Calculus 1}{Sep 2022 – Jan 2023}{Politecnico di Torino, Turin, Italy}
\entry{Teaching assistant – Physics 1}{Mar 2021 – May 2021}{Università degli Studi di Torino, Turin, Italy}

%==============================================================================
\section{Schools and conferences}
\begin{itemize}
  \item 50th Summer School on Mathematical Physics, Ravello, Italy (Aug 2025)
  \item Boltzmann Equation and Irreversibility: after 150 years an evergreen problem, Pisa, Italy (Nov 2022)
  \item XLVII Summer School on Mathematical Physics, Ravello, Italy (Aug 2022)
  \item Machine Learning Crash Course 2022, Genoa, Italy (Jun 2022)
\end{itemize}

%==============================================================================
\section{Skills}
\subsection{Computational}
\begin{itemize}
  \item \textbf{C++}: daily use, for numerical methods and code development.
  \item \textbf{Python}: daily use, for data analysis, scientific computing and interfaces to C/C++ codes.
  \item \textbf{MATLAB / Octave}: daily use, for numerical methods and code development.
  \item \textbf{Bash}: basic knowledge, for scripting.
  \item \textbf{Fortran}: basic knowledge, mainly for reading code.
\end{itemize}
\subsection{Operating systems and documents}
\begin{itemize}
  \item \textbf{Linux}, \textbf{macOS}: daily use.
  \item \textbf{LaTeX}: daily use. \textbf{Microsoft Office}: intermediate.
\end{itemize}
\subsection{Languages}
\begin{itemize}
  \item \textbf{Italian}: native. \textbf{English}: fluent.
\end{itemize}

%==============================================================================
\section{Authorization to personal data processing}
{\small\color{muted}
In compliance with the GDPR and Italian Legislative Decree no.~196 of 30/06/2003, I authorize the recipient of this document to use and process my personal details for the purpose of recruiting and selecting staff, and I confirm that I am informed of my rights under art.~7 of the above-mentioned Decree.

\smallskip
\begin{otherlanguage}{italian}
Autorizzo il trattamento dei dati personali ai sensi del GDPR 2016/679 del 27.04.2016 (Regolamento europeo relativo alla protezione delle persone fisiche con riguardo al trattamento dei dati personali).

\smallskip
Autorizzo la pubblicazione del Curriculum Vitae sul sito istituzionale del Politecnico di Milano (sez. Amministrazione Trasparente) in ottemperanza al D.~Lgs.~n.~33 del 14.03.2013 (e s.m.i.).\par
\end{otherlanguage}}

\end{document}
